This paper presented PostureCare, an edge vision IoT--LLM framework for continuous ergonomic risk
screening during computer use. The system combines a Freenove ESP32-S3 with an OV3660 camera, an
ultrasonic rangefinder and a photoresistor, and fuses head orientation, blink rate, viewing
distance, ambient illuminance and exposure duration into a single ergonomic state on low-cost,
non-contact hardware.

Its distinguishing element is what happens after detection. Rather than acting on instantaneous
threshold crossings, the system qualifies each detection by how long it persists, classifies the
qualified set into a four-tier severity model in which co-occurring mild conditions outrank isolated
ones, and passes the result to a notification governor that bounds how often the user is interrupted
through deduplication, exponential backoff, escalation with acknowledgement, and user-controlled
suppression. A cloud language model operates within those bounds as an explanation layer, and each
delivered reminder opens an observation window recording whether the condition cleared.

The contribution is an architecture and its implementation. All components described here are
built, and the design choices that constrain deployment --- a single configuration source, a single
writer for user-facing state, and a deterministic fallback for the recommendation layer --- are
recorded in Section~\ref{sec:implementation}. The system has not been evaluated: detection
accuracy, sensor calibration, interruption load and behavioural effectiveness are all unmeasured,
and a full protocol for establishing them is specified in Section~\ref{sec:protocol}.

Executing that protocol is the immediate next step. Beyond it, three directions follow from the
present design: separating severity per ergonomic dimension so that postural and visual conditions
can be reasoned about independently; adding daily cumulative exposure so that the temporal variable
connects to the epidemiological literature this work currently holds at arm's length; and moving
inference on-device, which would remove the privacy cost the current architecture accepts in
exchange for a richer landmark model.
